%============================================================
% CH1 · PATTERN 2 — REACHABILITY: THEORY
%============================================================
\subsection{Pattern Theory}
\label{ssec:ch1-p2-theory}

% ---- one reusable drawing of the running example G (8 vertices, 10 directed edges) ----
% every vertex style g1p2vK and every edge style g1p2eXY can be overridden per picture
\tikzset{
    g1p2v0/.style={gv}, g1p2v1/.style={gv}, g1p2v2/.style={gv}, g1p2v3/.style={gv},
    g1p2v4/.style={gv}, g1p2v5/.style={gv}, g1p2v6/.style={gv}, g1p2v7/.style={gv},
    g1p2e/.style={gde},
    g1p2e01/.style={g1p2e}, g1p2e02/.style={g1p2e}, g1p2e13/.style={g1p2e}, g1p2e23/.style={g1p2e},
    g1p2e34/.style={g1p2e}, g1p2e41/.style={g1p2e}, g1p2e50/.style={g1p2e}, g1p2e56/.style={g1p2e},
    g1p2e67/.style={g1p2e}, g1p2e76/.style={g1p2e},
    g1p2tree/.style={gtree, -{Stealth[length=2.6mm]}},
    pics/g1p2G/.style={code={
        \node[g1p2v0] (-0) at (0,0)        {0};
        \node[g1p2v1] (-1) at (1.5,0.85)   {1};
        \node[g1p2v2] (-2) at (1.5,-0.85)  {2};
        \node[g1p2v3] (-3) at (3,0)        {3};
        \node[g1p2v4] (-4) at (3,1.7)      {4};
        \node[g1p2v5] (-5) at (-1.5,0)     {5};
        \node[g1p2v6] (-6) at (-1.5,-1.4)  {6};
        \node[g1p2v7] (-7) at (-3,-1.4)    {7};
        \draw[g1p2e01] (-0) -- (-1);
        \draw[g1p2e02] (-0) -- (-2);
        \draw[g1p2e13] (-1) -- (-3);
        \draw[g1p2e23] (-2) -- (-3);
        \draw[g1p2e34] (-3) -- (-4);
        \draw[g1p2e41] (-4) -- (-1);
        \draw[g1p2e50] (-5) -- (-0);
        \draw[g1p2e56] (-5) -- (-6);
        \draw[g1p2e67] ([yshift=1.5mm]-6.west) -- ([yshift=1.5mm]-7.east);
        \draw[g1p2e76] ([yshift=-1.5mm]-7.east) -- ([yshift=-1.5mm]-6.west);
    }}
}

\begin{patternbox}{Pattern 2 at a Glance}
\textbf{Question:} \emph{starting from $s$ and following edges, which vertices can I get to?} The answer is the \textbf{reach set} $R(s)$. One search --- DFS or BFS, it does not matter which --- marks exactly the vertices of $R(s)$, touching every vertex and every edge at most once.

\begin{center}
\begin{tikzpicture}[g1p2v0/.style={gvsrc}, g1p2v1/.style={gvdone}, g1p2v2/.style={gvdone},
    g1p2v3/.style={gvdone}, g1p2v4/.style={gvdone}, g1p2v5/.style={gvdead}, g1p2v6/.style={gvdead},
    g1p2v7/.style={gvdead}, g1p2e01/.style={g1p2tree}, g1p2e02/.style={g1p2tree},
    g1p2e23/.style={g1p2tree}, g1p2e34/.style={g1p2tree},
    g1p2e56/.style={gfaded, -{Stealth[length=2.4mm]}}, g1p2e67/.style={gfaded, -{Stealth[length=2.4mm]}},
    g1p2e76/.style={gfaded, -{Stealth[length=2.4mm]}}]
    \pic (G) {g1p2G};
    \node[draw=successgreen, dashed, rounded corners=6pt, line width=1pt, fit=(G-0)(G-1)(G-2)(G-3)(G-4),
          inner sep=6pt, label={[font=\small\bfseries, text=successgreen]right:$R(0)=\{0,1,2,3,4\}$}] {};
    \node[note, anchor=north] at (-2.25,-1.9) {5, 6, 7 are not reachable\\(5 reaches 0, not vice versa)};
\end{tikzpicture}
\end{center}

\begin{lstlisting}[style=snippet]
seen[s] = true;  todo = {s};               // todo = a stack (DFS) or a queue (BFS)
while (todo not empty) {
    u = take one vertex out of todo;
    for (v in adj[u])
        if (!seen[v]) { seen[v] = true; add v to todo; }   // mark when ADDED
}
// now seen[v] == true  <=>  v is reachable from s.      Time O(|V| + |E|).
\end{lstlisting}
The thick blue edges discovered a new vertex: they form a tree rooted at $s$, the \emph{search tree}.
\end{patternbox}

%------------------------------------------------------------
\subsubsection*{From Pattern 1 to Pattern 2: from local to global}

Pattern~1 answered questions by looking at \textbf{one vertex and the edges touching it}: a degree, an in-degree, a pair of degrees. Such questions are \emph{local} --- one pass over the edge list suffices, and no vertex ever needs to know about a vertex two steps away.

Reachability is the first \emph{global} question. Whether $t$ can be reached from $s$ depends on whole chains of edges, possibly very long ones, and a single missing link anywhere in a chain changes the answer. Degree counting cannot see this. The two graphs below have \emph{identical} degrees --- every vertex has degree~2 --- yet in the first one $0$ reaches $4$ and in the second one it does not.

\begin{diagrambox}
\begin{tikzpicture}
    \begin{scope}
        \foreach \i in {0,...,5} \node[gv] (a\i) at ({90-60*\i}:1.25) {\i};
        \node[gvsrc] at (a0) {0}; \node[gvdst] at (a4) {4};
        \foreach \i/\j in {0/1,1/2,2/3,3/4,4/5,5/0} \draw[ge] (a\i) -- (a\j);
        \draw[gpath] (a0) -- (a5) -- (a4);
        \node[gv] at (a5) {5}; \node[gvsrc] at (a0) {0}; \node[gvdst] at (a4) {4};
        \node[note] at (0,-2) {one hexagon: $0 \leadsto 4$ \yes};
    \end{scope}
    \begin{scope}[xshift=5.2cm]
        \foreach \i/\x/\y in {0/0/1.1, 1/0.95/-0.45, 2/-0.95/-0.45} \node[gv] (b\i) at (\x,\y) {\i};
        \foreach \i/\x/\y in {3/3.6/1.1, 4/4.55/-0.45, 5/2.65/-0.45} \node[gv] (b\i) at (\x,\y) {\i};
        \node[gvsrc] at (b0) {0}; \node[gvdst] at (b4) {4};
        \foreach \i/\j in {0/1,1/2,2/0,3/4,4/5,5/3} \draw[ge] (b\i) -- (b\j);
        \draw[gcut] (1.8,1.5) -- (1.8,-1.0);
        \node[note] at (1.8,-1.55) {two triangles: $0 \not\leadsto 4$ \no};
    \end{scope}
\end{tikzpicture}
\figcap{Same number of vertices, same number of edges, same degree at every vertex --- different answers. A \emph{local} statistic cannot tell whether a chain is broken somewhere; a search can.}
\end{diagrambox}

So we need a procedure that \emph{walks}. Lesson~\ref{ssec:ch1-first-traversal} introduced the two classical walkers, DFS and BFS. In this pattern we use them for their most basic job: \textbf{deciding who can reach whom}. Their deeper properties (DFS edge types and timestamps, BFS distances) wait for Chapter~2 (DFS) and Chapter~3 (BFS).

%------------------------------------------------------------
\subsubsection*{Direction matters}

In an \textbf{undirected} graph reachability is symmetric: if $s$ reaches $t$, then $t$ reaches $s$ along the same path walked backwards. The vertices then split into groups that all reach each other --- the \emph{connected components} (Chapter~4). In a \textbf{directed} graph this symmetry is gone:

\begin{diagrambox}
\begin{tikzpicture}[x=1.4cm]
    \begin{scope}
        \node[gvsrc] (u) at (0,0) {u}; \node[gv] (v) at (1,0) {v}; \node[gv] (w) at (2,0) {w};
        \draw[ge] (u) -- (v); \draw[ge] (v) -- (w);
        \node[note] at (1,-0.8) {undirected: $u - v - w$,\\everybody reaches everybody};
    \end{scope}
    \begin{scope}[xshift=4.2cm]
        \node[gvsrc] (u) at (0,0) {u}; \node[gv] (v) at (1,0) {v}; \node[gv] (w) at (2,0) {w};
        \draw[gde] (u) -- (v); \draw[gde] (v) -- (w);
        \node[note] at (1,-0.8) {directed: $u$ reaches $w$,\\but $w$ reaches nobody else};
    \end{scope}
    \begin{scope}[xshift=8.4cm]
        \node[gv] (u) at (0,0) {u}; \node[gv] (v) at (1,0) {v}; \node[gv] (w) at (2,0) {w};
        \draw[gde] (u) -- (v); \draw[gde] (v) -- (w);
        \draw[gde] (w) to[bend right=50] (u);
        \node[note] at (1,-0.8) {directed cycle: symmetric again\\(strongly connected, Chapter 19)};
    \end{scope}
\end{tikzpicture}
\figcap{The same three vertices, three different reachability relations.}
\end{diagrambox}

Two consequences you will use constantly:
\begin{itemize}[leftmargin=*]
    \item An undirected edge $\{u,v\}$ is stored as \textbf{two} directed edges $u\to v$ and $v \to u$; then the same search code serves both kinds of graphs.
    \item In a directed graph, ``\emph{can $s$ reach $t$?}'' and ``\emph{can $t$ reach $s$?}'' are different questions; the second one is answered by searching from $t$ in the \textbf{reverse graph} (every edge flipped) --- see below.
\end{itemize}

%------------------------------------------------------------
\subsubsection*{When to recognise it}

\textbf{Use this pattern when the statement asks}
\begin{itemize}[leftmargin=*]
    \item ``is there a path / a way / a sequence of moves from $A$ to $B$?'' (yes / no);
    \item ``can we visit / open / infect / activate \emph{all} of them?'' --- is the reach set everything?;
    \item ``how many things are affected if this one starts?'' --- the size of a reach set (chain reactions, spreading rumours, detonations);
    \item ``starting from \emph{any} of these, what is reached?'' --- one multi-source search;
    \item ``who can reach $t$?'' --- one search in the reverse graph;
    \item ``for every pair \ldots'' with $|V|$ at most a few thousand --- one search per start vertex.
\end{itemize}

\textbf{Do NOT use a plain reachability search when}
\begin{itemize}[leftmargin=*]
    \item the question is \emph{how far} (fewest moves) --- that is BFS distances (Pattern~4, Chapter~3) or Dijkstra with weights (Chapter~7);
    \item the question only involves the edges of one vertex --- degrees suffice (Pattern~1) and are much simpler;
    \item there are many queries ``are $a$ and $b$ connected?'' in an undirected graph --- label the components once (Chapter~4) or use DSU, especially if edges keep arriving (Chapter~6);
    \item there are many queries ``does $a$ reach $b$?'' in a \emph{large} directed graph --- no simple trick exists in general; on DAGs use bitsets (Chapter~5), in general first compress the cycles (Chapter~19);
    \item every vertex has exactly one outgoing edge --- you are not searching, you are \emph{following arrows} (Pattern~5).
\end{itemize}

%------------------------------------------------------------
\subsubsection*{The algorithm and why it works}

Write $s \leadsto t$ for ``there is a walk from $s$ to $t$'' (a walk of zero edges is allowed, so $s \leadsto s$ always). The reach set is $R(s) = \{ t : s \leadsto t\}$.

The generic search keeps several kinds of vertices, shown in every trace of this book with the same colours:
\begin{center}\small
\legendchip{gvsrc}{start} \quad \legendchip{gv}{unseen} \quad \legendchip{gvseen}{marked, still in the to-do list} \quad \legendchip{gvcur}{being processed} \quad \legendchip{gvdone}{finished} \quad \legendchip{gvdead}{never reached}
\end{center}

Here is the iterative DFS (a stack) on the running example $G$, starting at $0$, with the adjacency lists in input order: $\Adj{0}=(1,2)$, $\Adj{1}=(3)$, $\Adj{2}=(3)$, $\Adj{3}=(4)$, $\Adj{4}=(1)$.

\begin{table}[H]
\centering\small
\begin{tabular}{@{}clll@{}}
\toprule
\textbf{step} & \textbf{pop} & \textbf{newly marked (pushed)} & \textbf{stack after the step (top on the right)} \\
\midrule
1 & 0 & 1, 2 & 1 2 \\
2 & 2 & 3    & 1 3 \\
3 & 3 & 4    & 1 4 \\
4 & 4 & --- \ (1 is already marked) & 1 \\
5 & 1 & --- \ (3 is already marked) & (empty) \\
\bottomrule
\end{tabular}
\end{table}

\begin{diagrambox}
\begin{tikzpicture}[scale=0.7, transform shape]
    \begin{scope}[g1p2v0/.style={gvdone}, g1p2v1/.style={gvseen}, g1p2v2/.style={gvdone},
        g1p2v3/.style={gvcur}, g1p2v4/.style={gvseen},
        g1p2e01/.style={g1p2tree}, g1p2e02/.style={g1p2tree},
        g1p2e23/.style={g1p2tree}, g1p2e34/.style={g1p2tree}]
        \pic (A) {g1p2G};
        \node[font=\large\bfseries] at (0,2.6) {after step 3: stack $= [1, 4]$};
    \end{scope}
    \begin{scope}[xshift=9cm, g1p2v0/.style={gvsrc}, g1p2v1/.style={gvdone}, g1p2v2/.style={gvdone},
        g1p2v3/.style={gvdone}, g1p2v4/.style={gvdone},
        g1p2v5/.style={gvdead}, g1p2v6/.style={gvdead}, g1p2v7/.style={gvdead},
        g1p2e01/.style={g1p2tree}, g1p2e02/.style={g1p2tree},
        g1p2e23/.style={g1p2tree}, g1p2e34/.style={g1p2tree}]
        \pic (B) {g1p2G};
        \node[font=\large\bfseries] at (0,2.6) {after step 5: stack empty};
    \end{scope}
\end{tikzpicture}
\figcap{Left: 0 and 2 are finished, 3 was just processed (it pushed 4), 1 and 4 wait on the stack. Right: the final state; the thick edges $0{\to}1$, $0{\to}2$, $2{\to}3$, $3{\to}4$ form the search tree.}
\end{diagrambox}

\begin{invariantbox}[Why the marked set is exactly the reach set]
\textbf{Invariant (soundness).} Every marked vertex is reachable from $s$.\\
\emph{Proof.} $s$ is marked first and $s \leadsto s$. Any other vertex $v$ is marked while scanning an edge $u \to v$ of an already marked vertex $u$; by induction $s \leadsto u$, hence $s \leadsto v$.

\smallskip
\textbf{Closure at the end (completeness).} When the to-do list is empty, \emph{no edge leads from a marked vertex to an unmarked one}.\\
\emph{Proof.} Every marked vertex was added to the to-do list exactly once, so it was taken out and scanned exactly once; during that scan every out-neighbour that was not yet marked got marked.

\smallskip
\textbf{Conclusion.} Let $s = v_0 \to v_1 \to \dots \to v_k = t$ be any walk. If $t$ were unmarked, take the first unmarked vertex $v_j$ on the walk ($j \ge 1$ because $v_0 = s$ is marked); then $v_{j-1} \to v_j$ is an edge from a marked to an unmarked vertex, contradicting closure. So every reachable vertex is marked, and by soundness only reachable vertices are. \hfill$\blacksquare$

\smallskip
\textbf{Running time.} A vertex is marked once, so it enters the to-do list once and its adjacency list is scanned once: $\sum_u (1 + \outdeg{u}) = |V| + |E|$ steps for a directed graph, $|V| + 2|E|$ for an undirected one --- $\VE$ either way. Memory: the graph, one \texttt{seen} flag per vertex and a to-do list of at most $|V|$ entries.
\end{invariantbox}

Notice what the proof did \emph{not} use: the order in which vertices leave the to-do list. Stack, queue, recursion, even a random bag --- all give the same marked set. The order matters only for the \emph{extra} information a search provides: DFS timestamps (Chapter~2) and BFS distances (Chapter~3). For pure reachability, \textbf{pick whatever you can write without bugs} --- in this book that is usually the iterative stack version.

The closure argument has a picture that is worth remembering:

\begin{diagrambox}
\begin{tikzpicture}[g1p2v0/.style={gvsrc}, g1p2v1/.style={gvdone}, g1p2v2/.style={gvdone},
    g1p2v3/.style={gvdone}, g1p2v4/.style={gvdone}, g1p2v5/.style={gvdead}, g1p2v6/.style={gvdead},
    g1p2v7/.style={gvdead}, g1p2e50/.style={gpath, -{Stealth[length=3mm]}}]
    \pic (G) {g1p2G};
    \draw[gcut] (-0.75,2.2) -- (-0.75,-2.0);
    \node[note, anchor=west, text width=4.6cm, align=left] at (3.9,0.6) {\textbf{inside the red line:} $R(0)$\\\textbf{outside:} 5, 6, 7\\[3pt]The only edge crossing the line is $5 \to 0$, and it points \textbf{inward}. No edge leaves $R(0)$ --- that is exactly why the search stopped there.};
\end{tikzpicture}
\figcap{Any set containing $s$ with no edge leaving it contains all of $R(s)$, and the search returns the \emph{smallest} such set. ``Here is a set with no exit'' is how you \emph{prove} that something is impossible in a reachability problem.}
\end{diagrambox}

%------------------------------------------------------------
\subsubsection*{Two tools that come for free: many sources and the reverse graph}

\textbf{Many sources.} ``What can be reached from \emph{at least one} of $s_1, \dots, s_k$?'' Do \emph{not} run $k$ searches. Mark all sources at the start, put them all into the to-do list and run the same loop once: $\VE$ in total. (Imagine a super-source $\star$ with an edge to every $s_i$; the multi-source search is a single search from $\star$.)

\textbf{The reverse graph.} ``Which vertices can reach $t$?'' Flip every edge: $u \to v$ becomes $v \to u$. A walk $u \leadsto t$ in $G$ is exactly a walk $t \leadsto u$ read backwards in the reverse graph $G^{R}$, so one search from $t$ in $G^{R}$ answers the question.

\begin{diagrambox}
\begin{tikzpicture}[scale=0.7, transform shape]
    \begin{scope}[g1p2v6/.style={gvsrc}, g1p2v2/.style={gvsrc}, g1p2v3/.style={gvdone},
        g1p2v4/.style={gvdone}, g1p2v1/.style={gvdone}, g1p2v7/.style={gvdone}]
        \pic (A) {g1p2G};
        \node[font=\large\bfseries, align=center] at (0,2.8) {multi-source search from $\{6, 2\}$\\reaches $\{1,2,3,4,6,7\}$};
    \end{scope}
    \begin{scope}[xshift=9cm, g1p2e/.style={ge, {Stealth[length=2.4mm]}-},
        g1p2v3/.style={gvsrc}, g1p2v0/.style={gvdone}, g1p2v1/.style={gvdone}, g1p2v2/.style={gvdone},
        g1p2v4/.style={gvdone}, g1p2v5/.style={gvdone}, g1p2v6/.style={gvdead}, g1p2v7/.style={gvdead}]
        \pic (B) {g1p2G};
        \node[font=\large\bfseries, align=center] at (0,2.8) {reverse graph $G^R$, search from $3$:\\who can reach $3$? $\{0,1,2,3,4,5\}$};
    \end{scope}
\end{tikzpicture}
\figcap{Left: both sources start marked; 0 and 5 stay white because no edge enters them from the reached part. Right: every arrow of $G$ is flipped; 6 and 7 cannot reach 3 (they only reach each other).}
\end{diagrambox}

Combining the two tools answers ``\emph{which vertices lie on some walk from $s$ to $t$?}'': exactly those in $R(s) \cap R^{R}(t)$ --- forward from $s$ intersected with backward from $t$ (Exercise~E2.6).

%------------------------------------------------------------
\subsubsection*{Variations map}

\begin{diagrambox}
\begin{tikzpicture}[
    var/.style={rectangle, rounded corners=2pt, draw=patternborder, fill=patternbg,
                font=\scriptsize, align=center, text width=3.25cm, minimum height=1.25cm, inner sep=3pt}]
    \node[var, fill=answerfill, draw=accentorange, text width=4.6cm, font=\scriptsize\bfseries] (root) at (0,0)
        {Reachability from $s$\\one DFS/BFS, mark on push, $\VE$};
    \node[var] (v1) at (-5.4,-2) {\textbf{yes/no: $s \leadsto t$?}\\stop as soon as $t$ appears\\(Problem 1.2.1)};
    \node[var] (v2) at (-1.8,-2) {\textbf{``reach everything?''}\\number marked $= n$?\\(Problem 1.2.2)};
    \node[var] (v3) at (1.8,-2)  {\textbf{size of $R(s)$}\\chain reactions: count marks\\(E2.5)};
    \node[var] (v4) at (5.4,-2)  {\textbf{many sources}\\push them all first,\\one search};
    \node[var] (v5) at (-5.4,-3.6) {\textbf{who reaches $t$?}\\search the reverse graph\\(E2.4, E2.6)};
    \node[var] (v6) at (-1.8,-3.6) {\textbf{all pairs}, small $n$\\one search per vertex\\$\bigO{|V|(|V|+|E|)}$ (E2.3)};
    \node[var] (v7) at (1.8,-3.6)  {\textbf{edges given by a rule}\\keys, radii, portals:\\build adj.\ first (E2.1)};
    \node[var] (v8) at (5.4,-3.6)  {\textbf{many $a$--$b$ queries}\\undirected: components\\or DSU (Ch.~4, Ch.~6)};
    \node[draw=patternborder, dashed, rounded corners=4pt, fit=(v1)(v4)(v5)(v8), inner sep=5pt] (grp) {};
    \draw[flow] (root.south) -- (grp.north);
\end{tikzpicture}
\end{diagrambox}

\textbf{Twists and how to handle them}
\begin{itemize}[leftmargin=*]
    \item \textbf{Forbidden vertices or edges} (walls, closed roads, ``you may not pass through $x$''): never mark them, never follow them. Setting \texttt{seen[x] = 1} for every forbidden $x$ before the search starts does it in one line.
    \item \textbf{Edges defined by a rule} (a key lying in a room, a bomb's radius, a portal $i \to i + a_i$): either build the adjacency lists first or generate the neighbours inside the loop. For $|V| \le 100$ with an $\bigO{|V|^2}$ rule, building the full lists is simplest (E2.5).
    \item \textbf{Every vertex has one out-edge:} the ``search'' degenerates into a loop \texttt{cur = next[cur]} (E2.1, and all of Pattern~5).
    \item \textbf{Repeated queries from different sources:} each query costs a fresh search unless the graph is undirected (then label the components once). If $q \cdot n$ is too big to clear a \texttt{seen} array per query, a \emph{timestamp} array avoids the clearing (Mistake~5).
    \item \textbf{``Fewest starting points that reach everything'':} in a DAG these are exactly the vertices with in-degree~0 (E2.2) --- a reachability question answered by Pattern~1.
\end{itemize}

%------------------------------------------------------------
\subsubsection*{Code templates}

All four routines below are $\VE$; the all-pairs loop at the end of \texttt{main()} is $\bigO{|V|(|V|+|E|)}$. The \texttt{main()} uses the running example $G$; the values in the comments were printed by the program.

\begin{lstlisting}[style=cpp]
#include <bits/stdc++.h>
using namespace std;

// (1) iterative DFS: mark every vertex reachable from s (no recursion: safe for 10^6 vertices)
vector<char> reachFrom(const vector<vector<int>>& adj, int s) {
    vector<char> seen(adj.size(), 0);
    vector<int> st(1, s);
    seen[s] = 1;                                   // mark when PUSHED, never twice
    while (!st.empty()) {
        int u = st.back(); st.pop_back();
        for (int v : adj[u])
            if (!seen[v]) { seen[v] = 1; st.push_back(v); }
    }
    return seen;
}

// (2) BFS from a whole SET of sources at once (one search, not one per source)
vector<char> reachFromSet(const vector<vector<int>>& adj, const vector<int>& sources) {
    vector<char> seen(adj.size(), 0);
    queue<int> q;
    for (int s : sources)
        if (!seen[s]) { seen[s] = 1; q.push(s); }
    while (!q.empty()) {
        int u = q.front(); q.pop();
        for (int v : adj[u])
            if (!seen[v]) { seen[v] = 1; q.push(v); }
    }
    return seen;
}

// (3) the reverse graph: u -> v becomes v -> u.  "Who can reach t?" = reachFrom(rev, t)
vector<vector<int>> reversedGraph(const vector<vector<int>>& adj) {
    vector<vector<int>> rev(adj.size());
    for (int u = 0; u < (int)adj.size(); u++)
        for (int v : adj[u]) rev[v].push_back(u);
    return rev;
}

// (4) recursive DFS: shortest to write, but the depth can reach |V| (see Mistake 3)
void dfsRec(int u, const vector<vector<int>>& adj, vector<char>& seen) {
    seen[u] = 1;
    for (int v : adj[u])
        if (!seen[v]) dfsRec(v, adj, seen);
}

void show(const char* name, const vector<char>& seen) {
    cout << name << ":";
    for (int v = 0; v < (int)seen.size(); v++)
        if (seen[v]) cout << " " << v;
    cout << "\n";
}

int main() {
    int n = 8;
    int E[][2] = {{0,1},{0,2},{1,3},{2,3},{3,4},{4,1},{5,0},{5,6},{6,7},{7,6}};
    vector<vector<int>> adj(n);
    for (auto& e : E) adj[e[0]].push_back(e[1]);         // directed edges

    show("reach(0)", reachFrom(adj, 0));                  // reach(0): 0 1 2 3 4
    vector<char> seen(n, 0);
    dfsRec(0, adj, seen);
    show("dfsRec(0)", seen);                              // dfsRec(0): 0 1 2 3 4
    show("reach(6)", reachFrom(adj, 6));                  // reach(6): 6 7
    show("canReach(3)", reachFrom(reversedGraph(adj), 3)); // canReach(3): 0 1 2 3 4 5
    vector<int> src = {6, 2};
    show("reach({6,2})", reachFromSet(adj, src));         // reach({6,2}): 1 2 3 4 6 7

    long long pairs = 0;                                  // all pairs (s, t) with s ->* t
    for (int s = 0; s < n; s++) {
        vector<char> r = reachFrom(adj, s);
        pairs += count(r.begin(), r.end(), 1);
    }
    cout << "reachable pairs = " << pairs << "\n";        // reachable pairs = 30
}
\end{lstlisting}

\begin{insightbox}[Check the templates by hand]
\begin{itemize}[leftmargin=*]
    \item \textbf{reach(0)} $=\{0,1,2,3,4\}$ (the trace above); \textbf{reach(6)} $= \{6, 7\}$: the 2-cycle $6 \rightleftarrows 7$ has no exit.
    \item \textbf{canReach(3)}: in $G^R$ the search from 3 follows $3 \to 1$ and $3 \to 2$ (from $1{\to}3$, $2{\to}3$), then $1 \to 0$, $1 \to 4$, $2 \to 0$, $0 \to 5$, and stops: $\{0,1,2,3,4,5\}$.
    \item \textbf{All pairs:} $|R(v)|$ for $v = 0, \dots, 7$ is $5, 3, 4, 3, 3, 8, 2, 2$, and the sum is $30$ (pairs $(v,v)$ included).
\end{itemize}
\end{insightbox}

\begin{tricksbox}[Which version should I type?]
\begin{itemize}[leftmargin=*]
    \item \textbf{Yes/no or ``how many are reached'':} the iterative stack version (1). It is as short as the recursive one and never overflows.
    \item \textbf{Recursive DFS (4)} is fine when the depth is provably small: $|V|$ up to about $10^4$ with a default 1--8\,MB stack, or a tree of known small height. A path graph with $2\cdot 10^5$ vertices makes the recursion $2\cdot10^5$ frames deep --- a likely \stackoverflow{} on many judges.
    \item \textbf{BFS (2)} costs the same and gives distances for free --- use it whenever the next question might be ``in how many steps?''.
    \item \textbf{\texttt{vector<char>}} instead of \texttt{vector<bool>} for \texttt{seen}: one byte per vertex, no bit fiddling, usually a little faster.
\end{itemize}
\end{tricksbox}

%------------------------------------------------------------
\subsubsection*{Common mistakes}

\begin{warningbox}[Mistake 1 --- marking when popped instead of when pushed]
\begin{lstlisting}[style=snippet]
// WRONG: v may be pushed once per incoming edge -> the stack grows to O(|E|),
//        and every copy re-scans adj[v] unless you also check on pop
for (int v : adj[u]) if (!seen[v]) st.push_back(v);      // seen[v] set only when popped
// RIGHT: mark at the moment of discovery; each vertex enters the stack once
for (int v : adj[u]) if (!seen[v]) { seen[v] = 1; st.push_back(v); }
\end{lstlisting}
``Mark on pop'' is correct only with an extra \texttt{if (seen[u]) continue;} right after popping, and it still needs $\bigO{|E|}$ memory. Marking on push is simpler and never worse.
\end{warningbox}

\begin{warningbox}[Mistake 2 --- storing an undirected edge once]
\begin{lstlisting}[style=snippet]
adj[a].push_back(b);                          // WRONG if undirected: b cannot go back
adj[a].push_back(b); adj[b].push_back(a);     // RIGHT
\end{lstlisting}
Reread the statement: ``bi-directional'', ``roads can be used in both directions'' $\Rightarrow$ two entries; ``one-way'', ``from $a$ to $b$'', keys, portals $\Rightarrow$ one entry.
\end{warningbox}

\begin{warningbox}[Mistake 3 --- deep recursion]
Recursive DFS on a path of $n = 2\cdot10^5$ vertices nests $2\cdot 10^5$ calls. Each frame costs tens of bytes or more, and the default stack is often 1\,MB (Windows) or 8\,MB (Linux): the result is a crash without a useful message. Use the stack version whenever $|V|$ can exceed about $10^4$.
\end{warningbox}

\begin{warningbox}[{Mistake 4 --- searching in the wrong direction}]
``Can every city reach the capital?'' is \emph{not} answered by a search \emph{from} the capital in a directed graph. Search from the capital in the \textbf{reverse} graph (or flip each edge while reading the input).
\end{warningbox}

\begin{warningbox}[Mistake 5 --- resetting \texttt{seen} the expensive way]
With many test cases or many queries, either re-create \texttt{seen} for each test (\texttt{seen.assign(n, 0)} costs $\bigO{n}$ of the \emph{current} test --- fine) or use a timestamp: \texttt{stamp[v] == query} means ``seen in this query'', and starting a new query is just \texttt{query++}. Never \texttt{memset} a $10^6$-sized array $10^5$ times.
\end{warningbox}

\begin{warningbox}[Mistake 6 --- 1-indexed input in a 0-indexed array]
Vertices numbered $1..n$ in the input: either allocate $n+1$ lists or subtract 1 when reading. Mixing the two silently drops vertex $n$ or creates a phantom vertex $0$ that is ``never reached''.
\end{warningbox}

\begin{checklistbox}[``I'm stuck'' checklist for Pattern 2]
\begin{itemize}[leftmargin=*]
    \item[$\square$] What is a vertex? What is an edge? \textbf{Directed or not?} (Write the Graph Model box.)
    \item[$\square$] Is the question really ``reach'' --- or ``how far'' (BFS distances), ``how many groups'' (components), ``one arrow per vertex'' (Pattern~5)?
    \item[$\square$] From where do I search? One source, many sources (push them all), or from the target in the reverse graph?
    \item[$\square$] Is the answer a yes/no, the number of marked vertices, or ``all $n$ marked''?
    \item[$\square$] Can I afford one search per vertex ($|V| \cdot (|V|+|E|)$ operations)?
    \item[$\square$] Is the recursion depth safe? Are both directions of undirected edges stored? Is \texttt{seen} reset between tests?
    \item[$\square$] To prove impossibility: can I exhibit a set containing $s$, not containing $t$, with no edge leaving it?
\end{itemize}
\end{checklistbox}
